\slajdid{02-s004}
\begin{frame}[label=02-s004]{Srednje kašnjenje i kratak impuls}
\setlength{\abovedisplayskip}{4pt}\setlength{\belowdisplayskip}{4pt}
% element: 02-s004:e01
\Rezultat{Zapamtiti: \(\textcolor{DEcrvena}{t_{pHL}\ne t_{pLH}}\).}
\par\vspace{2mm}
\begin{columns}[T,onlytextwidth]
\begin{column}{.67\textwidth}
% element: 02-s004:e02
Srednje kašnjenje logičkog kola:
\[t_p=\frac{t_{pHL}+t_{pLH}}{2}\]
Za pojednostavljen prikaz \alert{nekih} pojava uzimamo
\[t_{pHL}=t_{pLH}=t_p.\]
% element: 02-s004:e04
\textbf{Da li je izlazni signal samo zakašnjena verzija ulaznog?}\par\vspace{1mm}
Može se promeniti oblik i trajanje impulsa, a impuls može i da izostane.\par\vspace{1mm}
\alert{Ishod zavisi od trajanja ulaznog impulsa u odnosu na kašnjenja kola.}
\end{column}
\begin{column}{.29\textwidth}\centering
% element: 02-s004:e03
\Oznaka{Kratak impuls}\par\vspace{3mm}
\begin{tikzpicture}[DEoznaka,x=8.2mm,y=7.5mm]
\draw[DEstrelica] (-.7,2.25)--(3.8,2.25) node[below] {\(t\)};
\draw[DEstrelica] (0,2.10)--(0,4.05) node[left] {\(x(t)\)};
\draw[DEstrelica] (-.7,-.1)--(3.8,-.1) node[below] {\(t\)};
\draw[DEstrelica] (0,-.25)--(0,1.7) node[left] {\(y(t)\)};
\node[above left,inner sep=1pt] at (0,3.35) {=1};\node[above left,inner sep=1pt] at (0,2.45) {=0};
\node[above left,inner sep=1pt] at (0,1) {=1};\node[above left,inner sep=1pt] at (0,.1) {=0};
\draw[DEvod] (-.5,2.45)--(.6,2.45)--(.6,3.35)--(.85,3.35)--(.85,2.45)--(3.5,2.45);
\draw[DEvod] (-.5,1)--(2.8,1)--(2.8,.1)--(3.05,.1)--(3.05,1)--(3.5,1);
\draw[dashed] (.6,3.35)--(.6,-.35) (2.8,.1)--(2.8,-.35) (0,.1)--(2.8,.1);
\draw[DEcrvena,<->] (.6,-.3)--(2.8,-.3) node[midway,below] {\(t_p\)};
\end{tikzpicture}

\end{column}
\end{columns}
\end{frame}
